%% ============================================================
\section{Conclusion}
\label{sec:conclusion}
%% ============================================================

Autonomous research systems have reached the point where solver quality alone no longer differentiates them---multiple systems achieve competitive scores on the same benchmarks with vastly different approaches.
What separates their outputs is whether the resulting paper can be trusted.
Our 75-paper audit shows that no baseline produces papers free of evidence chain failures, and the failures---hallucinated references up to 21\%, fictional method sections, scores on the wrong scale---are undetectable by evaluations that only assess surface presentation rather than evidence grounding.

Chain-of-Evidence reframes verifiability as a first-class design constraint.
\sys{} demonstrates that an end-to-end pipeline can maintain evidence chains without sacrificing solver competitiveness, and \coeaudit{} provides a reusable procedure for auditing any system's output.
The gap between \sys{}'s results and the baselines confirms that verifiability is architectural: systems that build evidence chains at claim-production time produce more verifiable outputs than those that reconstruct grounding after the fact.
The harder problems---verifying citation support, checking conclusion claims, extending to domains without deterministic evaluators---remain open, but they are tractable extensions of the checks demonstrated here, and their importance grows with the volume of AI-generated research.
